\section{Teorien i praksis: fem utarbeidde døme}
\label{sec:examples-nn}

Den empiriske demonstrasjonen i seksjon~\ref{sec:empirical-nn} testar teoriens predikasjonar mot eit historisk korpus. Men ein teori om designform må òg kaste lys over samtidsdesign. Denne seksjonen anvender dei sju proposisjonane på fem konkrete tilfelle: eitt historisk, to samtidige og to framveksande. Døma er \textit{illustrative}, ikkje probative; dei demonstrerer teoriens forklaringskraft og genererer testbare predikasjonar, men dei utgjer ikkje uavhengige empiriske testar.

\subsection{Utkragingsstolkappløpet, 1925--1928}

Mellom 1925 og 1928 konvergerte fire designarar i tre byar uavhengig av kvarandre mot ein delt morfologisk innovasjon: den utkraga røyrstolstolen. Marcel Breuer designa Wassily (B3) i 1925 ved Bauhaus i Dessau, med kaldbøygd nikkelbelagt stålrøyr inspirert av Adler-sykkelen sin. Mart Stam presenterte den fyrste ekte utkraginga (S33) ved ei bustadmesse i Stuttgart i 1926, med rette gassrøyrsegment samanføydde i rette vinklar. Ludwig Mies van der Rohe stilte ut sin utkraga MR-stol i 1927 på Weissenhof-utstillinga, med eit samanhengande bøygd røyr som gav setet ein karakteristisk fjøring. Breuer returnerte med Cesca (B32) i 1928, som kombinerte ein utkraga stålramme med eit sete og rygg i rotting.

I teoriens termar hadde materialaffordansen til kaldbøygd stålrøyr vore latent i tiår (proposisjon 2). Stålrøyr vart produsert for syklar og røyrlegging frå 1880-åra. \textit{Operasjonen}; kaldbøying av tynnvegga røyr til møbelstorleikar; var den manglande lenka. Konvergensen av fire designarar innanfor tre år stadfestar at landskapsattraktoren var reell (proposisjon 1): fitnesstopen eksisterte i landskapet før nokon klattra han.

\subsection{Neri Oxman sitt Silk Pavilion, MIT Media Lab, 2013}

Silk Pavilion er ein kupelforma struktur produsert gjennom ein tofaseprosess. I den fyrste fasen deponerte ei CNC-styrt maskin eit stillas av silketrådar over ein stålramme. I den andre fasen vart 6\,500 silkeormar plasserte på stillaset og deponerte tilleggssilke over 20 dagar, styrte av miljøparametrar inkludert lysintensitet og temperaturgradi\-entar.

Resultatet er ein struktur der den endelege forma ikkje vart teikna av designaren. Neri Oxman og teamet hennar sette grensevilkåra; dei biologiske agentane navigerte formgrammatikken på cellulært nivå. Kvar silkeorm er ein null-ordens agent i teoriens forstand (proposisjon 5): han har ein måltilstand (spinne ein kokong), eit avstandsmål (nærleik til eigna festepunkt) og ein justeringsmekanisme (silkeekstruderingsretning). Ormen representerer ikkje kupelen; han responderer på den lokale gradienten.

\subsection{Figma Weave og designaren som landskapsnavigator}

I 2025 annonserte Figma Weave, ein generativ AI-funksjon innebygd i designverktøyet. Weave produserer grensesnittoppsett, komponentarrangeringar og visuelle designar frå tekstlege eller kontekstuelle prompt. Designaren gjennomgår, vel, modifiserer og kombinerer genererte output snarare enn å konstruere kvart element manuelt.

I teoriens termar vert formgrammatikken $SG$ no delvis utført av ein maskinagent. Designaren si rolle skiftar frå å \textit{utføre} grammatikkreglar til å \textit{styre den kognitive lyskjegla}: velje kva dimensjonar av morforommet systemet bør gje merksemd til. Kvaliteten på det resulterande designet avheng ikkje av AI-en sin generative kapasitet, men av dimensjonaliteten til mennesket si omsorg.

\subsection{Samtaleagentdesign}

Når ein interaksjonsdesignar designar ein samtaleagent; ein kundeservice-chatbot, ein medisinsk triageassistent, eit språktutoringssystem; designar ho ein navigasjonsalgoritme for ein ikkje-menneskeleg agent. Chatboten er ein agent $A := (g, d, \delta)$ med eit mål, eit avstandsmål og ein justeringsmekanisme.

Interaksjonsdesignaren er ein \textbf{meta-agent} der den kognitive lyskjegla må romme både chatboten si lyskjegle og brukaren si lyskjegle (proposisjon 6). Kvaliteten på interaksjonen avheng av kor mange aksar designaren aktivt representerer: oppgåveeffektivitet, brukarfrustrasjon, emosjonell tone, tilgjengelegheit, kulturell eignaheit. Kvar urepresentert akse er ein akse der designet vil svikte (proposisjon 5).

\subsection{The Great Green Wall: kontinental-skala økologisk design}

Great Green Wall-initiativet, lansert av Den afrikanske unionen i 2007, har som mål å restaurere 100 millionar hektar degradert land langs eit 8\,000 kilometer langt belte frå Senegal til Djibouti. Prosjektet adresserer ørkenspreiing, matusikkerheit, klimaendring og tvungen migrasjon over Sahel.

Dette er design i ein skala som direkte motseier Le Corbusier sitt prinsipp om mennesket som mål for alt. ``Forma'' under design er ikkje eit enkelt objekt, men ein multiskalar konfigurasjon: jordmikrobiom-samansetjing (cellenivå-agent), rotarkitektur og mykorrhiza-nettverk (organismenivå-agent), kronedekke og artsmangfald (samfunnsnivå-agent), pastoral arealbruk og samfunnsstyring (sosialt nivå-agent), kontinental nedbørsdynamikk og karbonsirkulasjon (systemnivå-agent).

Teorien predikerer (proposisjon 4) at å kollapse landskapet til eitt seleksjonstrykk; til dømes karbonbinding; vil produsere konfigurasjonar som sviktar under trykk designet var blindt for. Dei mest vellukka restaureringsstadene; som bøndestyrt naturleg regenerering i Niger; nytta ein flertrykk-tilnærming som koordinerte biologiske, sosiale og økonomiske agentar samstundes.
